\documentclass[a4paper]{article}
% Español (México) con XeLaTeX: fontspec para fuentes Unicode, polyglossia
% para la división silábica y los nombres del documento en la variante mexicana.
\usepackage{fontspec}
\usepackage{polyglossia}
\setdefaultlanguage[variant=mexican]{spanish}
% Los nombres chinos van como «pinyin (original)»: xeCJK da a los caracteres
% Han su propia fuente; sin ella XeLaTeX los omite con un aviso.
% FandolSong no cubre algunos caracteres de nombres (U+73FA); WenQuanYi Zen Hei sí.
\usepackage[AutoFallBack=true]{xeCJK}
\setCJKmainfont{FandolSong-Regular.otf}
\setCJKfallbackfamilyfont{\CJKrmdefault}{WenQuanYi Zen Hei}
% polyglossia no traduce los nombres de listings.
\AtBeginDocument{\providecommand{\lstlistingname}{}\renewcommand{\lstlistingname}{Listado}%
  \providecommand{\lstlistlistingname}{}\renewcommand{\lstlistlistingname}{Índice de listados}}
\usepackage{amsmath, amssymb}
\usepackage{graphicx}
\usepackage[margin=2.35cm]{geometry}
\usepackage[most]{tcolorbox}
\usepackage{booktabs}
\usepackage{tabularx}
\usepackage{longtable}
\usepackage{array}
\usepackage{float}
\usepackage{xcolor}
\usepackage{hyperref}
\usepackage{fancyhdr}
\setCJKmainfont[AutoFakeBold=2.4]{AR PL UMing CN}
\setCJKsansfont[AutoFakeBold=2.4]{Droid Sans Fallback}
\setCJKmonofont[AutoFakeBold=2.4]{Droid Sans Fallback}
\hypersetup{colorlinks=true, linkcolor=blue!55!black, urlcolor=blue!60!black}
\graphicspath{{./}{figures/}}
\setlength{\parskip}{0.55em}
\setlength{\parindent}{2em}
\setlength{\emergencystretch}{3em}
\renewcommand{\arraystretch}{1.18}
\newtcolorbox{knowledgebox}[1]{enhanced, colback=blue!5!white, colframe=blue!70!black, colbacktitle=blue!70!black, coltitle=white, fonttitle=\bfseries, title={#1}, sharp corners, breakable}
\newtcolorbox{importantbox}[1]{enhanced, colback=yellow!10!white, colframe=yellow!75!black, colbacktitle=yellow!75!black, coltitle=black, fonttitle=\bfseries, title={#1}, sharp corners, breakable}
\newtcolorbox{warningbox}[1]{enhanced, colback=red!5!white, colframe=red!70!black, colbacktitle=red!70!black, coltitle=white, fonttitle=\bfseries, title={#1}, sharp corners, breakable}
\pagestyle{fancy}
\fancyhf{}
\fancyfoot[C]{\thepage}
\fancyfoot[R]{\footnotesize\color{gray}\href{https://github.com/hqhq1025/ai-course-notes}{AI Course Notes}}
\renewcommand{\headrulewidth}{0pt}
\renewcommand{\footrulewidth}{0pt}
\newcommand{\videourl}{https://www.youtube.com/watch?v=ouG6jrkECrc}
\newcommand{\repourl}{https://github.com/hqhq1025/ai-course-notes}

\begin{document}
\begin{titlepage}
\centering
{\Large Notas didácticas de una entrevista a fondo\par}
\vspace{1.0cm}
{\huge\bfseries Kimi K2, Agentic LLM y estar en el comienzo del infinito}\par
\vspace{0.5cm}
{\Large Zhang Xiaojun conversa con Yang Zhilin (杨植麟), fundador de Moonshot AI (月之暗面)\par}
\vspace{0.35cm}
{\large Elaborado a partir del video público de Zhang Xiaojun Podcast, la transcripción hecha con GPU, la estructura de capítulos y diagramas conceptuales propios\par}
\vspace{0.3cm}
{\large 2025-08-27\par}
\vspace{0.85cm}
\includegraphics[width=0.88\textwidth,height=0.42\textheight,keepaspectratio]{cover.jpg}\par
\vfill
\begin{tcolorbox}[width=0.92\textwidth, colback=black!2!white, colframe=black!55, sharp corners]
\textbf{Título del video}: 113. Conversación con Yang Zhilin un año después: K2, Agentic LLM, el cerebro en una cubeta y «estar en el comienzo del infinito»\par
\textbf{Canal del video}: Zhang Xiaojun Podcast\par
\textbf{Duración del video}: 01:41:14\par
\textbf{Enlace del video}: \href{\videourl}{\nolinkurl{\videourl}}\par
\textbf{Notas sobre el material}: YouTube no tiene subtítulos disponibles; el cuerpo del texto se elaboró a partir de la transcripción con faster-whisper large-v3 en CUDA, los capítulos del video, la descripción del video, la portada y figuras didácticas propias. El video, de portada fija, no se reutiliza en el cuerpo del texto.\par
\textbf{Página del proyecto}: \href{\repourl}{\nolinkurl{\repourl}}\par
\end{tcolorbox}
\end{titlepage}

\tableofcontents
\newpage

\section{Introducción: por qué K2 es el «cerebro en una cubeta» al que le crecen manos}

Esta sección establece primero el hilo técnico de todo el episodio. Kimi K2 es un modelo de lenguaje de gran escala de código abierto, orientado a la programación y a lo Agentic, basado en la arquitectura MoE. La metáfora más contundente de la entrevista es la del «cerebro en una cubeta»: un modelo que solo sabe pensar y producir texto es como un cerebro colocado en una pecera, sin una conexión real con el mundo exterior; en cambio, cuando el modelo adquiere capacidad de programación, capacidad de usar herramientas en múltiples turnos y retroalimentación externa, empiezan a crecerle «manos» y puede manipular el mundo digital.

En este episodio hay que entender cuatro cosas. Primero, por qué K2 reúne el modelo base, la token efficiency, las capacidades Agentic y el código abierto. Segundo, por qué la test-time scaling apunta a la vez al RL de pensamiento largo y al Agent RL. Tercero, cómo se equilibran el muro de datos, el FLOPs scaling, la Linear Attention, las arquitecturas de contexto largo y el «coeficiente intelectual» del modelo. Cuarto, cómo Yang Zhilin (杨植麟) integra el juicio técnico, el modelo de negocio y la mentalidad de fundador en el relato de «estar al comienzo del infinito».

\begin{figure}[H]
\centering
\includegraphics[width=0.96\textwidth]{brain-in-vat-to-agent.png}
\caption{Del cerebro en una cubeta al Agent: la capacidad de programación y el uso de herramientas permiten que el modelo empiece a manipular el mundo digital. Diagrama conceptual propio, elaborado a partir de la conversación de 00:01:19--00:24:58.}
\end{figure}

\begin{knowledgebox}{Lectura de la figura: por qué las «manos» son decisivas}
Que un modelo sepa pensar y que sepa actuar son dos cosas distintas. El código, las llamadas a herramientas, las API, el sistema de archivos y los entornos de retroalimentación permiten que el razonamiento del modelo modifique el estado externo. Lo esencial de un Agentic LLM no es solo decir más pasos, sino usar herramientas en múltiples turnos y corregirse según los resultados.
\end{knowledgebox}

\begin{importantbox}{Tesis central del episodio}
La importancia de K2 no es solo ser «otro modelo potente», sino integrar en un mismo sistema el modelo base, la capacidad de programación, el uso de herramientas, el ecosistema de código abierto y la test-time scaling: llevar al modelo de un sistema cognitivo cerrado a un agente digital capaz de actuar.
\end{importantbox}

\begin{knowledgebox}{Panorama del sistema en este episodio}
\begin{tabularx}{\textwidth}{p{0.20\textwidth}p{0.34\textwidth}X}
\toprule
Nivel & Pregunta central & Respuesta en la entrevista \\
\midrule
Nivel filosófico & ¿Por qué seguir escalando? & Los problemas son inevitables, pero los problemas pueden resolverse. \\
Nivel del modelo & ¿Cómo se convierte K2 en una base sólida? & MoE, token efficiency, Muon, Rephrase. \\
Nivel de acción & ¿Cómo le crecen manos al modelo? & Capacidad de programación, uso de herramientas en múltiples turnos, generalización Agentic. \\
Nivel del ecosistema & ¿Por qué código abierto? & Para que las capacidades del modelo base se validen externamente y el ecosistema las amplifique. \\
Nivel comercial & ¿Cómo se gana dinero y dónde están los límites? & API, productos Agent y una nueva delimitación entre empresas de modelos base y empresas de aplicaciones. \\
Nivel organizacional & ¿Cómo soporta el fundador la complejidad? & Gestionar con retroalimentación al estilo RL, pero con cuidado de no caer en el hack de las métricas. \\
\bottomrule
\end{tabularx}
\end{knowledgebox}

\begin{knowledgebox}{Términos clave: palabras clave del episodio}
\begin{tabularx}{\textwidth}{p{0.22\textwidth}p{0.32\textwidth}X}
\toprule
Término & Significado & Por qué importa \\
\midrule
MoE & Mixture of Experts: se enruta cada token o tarea a distintos expertos & Aumenta la capacidad de parámetros y, al mismo tiempo, controla el cómputo activado en cada paso. \\
Agentic LLM & Modelo de lenguaje que puede usar herramientas en múltiples turnos, observar la retroalimentación y completar tareas & Pasa de la conversación a la acción en el mundo digital. \\
Test-time scaling & Invertir más cómputo, más pensamiento o más llamadas a herramientas durante la inferencia & Permite que el modelo siga ampliando sus capacidades en la fase de prueba. \\
Token efficiency & La ganancia de capacidad que aporta cada token de entrenamiento & En la era del muro de datos, con los mismos datos el «cerebro» debe crecer más. \\
Muon & Línea de optimizadores mencionada en la entrevista & Puede mejorar la eficiencia del entrenamiento, pero plantea retos de estabilidad. \\
\bottomrule
\end{tabularx}
\end{knowledgebox}

\subsection{Resumen de la sección}

EP113 es una entrevista de juicio técnico posterior al lanzamiento de K2. Sitúa al Agentic LLM dentro de la estrategia de las empresas de modelos, el código abierto, la eficiencia del entrenamiento, las decisiones de arquitectura y la psicología del fundador, y conviene leerla junto con la revisión de arquitecturas de Attention de EP119 y el marco teórico de Agent de EP115.
\section{Una montaña infinita: los problemas son inevitables, pero tienen solución}

Esta sección parte de «The Beginning of Infinity». Yang Zhilin (杨植麟) dice que los problemas son inevitables, pero que tienen solución; quizá esta montaña nevada no tenga cima, pero él espera que nunca la tenga. La metáfora describe bien a una empresa de modelos: cada vez que mejora la capacidad del modelo se desbloquean nuevos escenarios y, al mismo tiempo, surgen nuevos problemas. El preentrenamiento, el RLHF, el pensamiento largo, las capacidades Agentic, el uso de herramientas y los sistemas de evaluación son como rutas a distintas alturas de la montaña nevada.

\begin{figure}[H]
\centering
\includegraphics[width=0.92\textwidth]{infinite-mountain-map.png}
\caption{Una montaña infinita: los problemas son inevitables, pero tienen solución; el avance de los modelos abre continuamente nuevos problemas. Diagrama conceptual propio, elaborado a partir de la conversación de 00:00:02--00:05:18.}
\end{figure}

\begin{knowledgebox}{Lectura de la figura: lo infinito no es pesimismo, sino el oxígeno de la investigación}
Si la montaña tuviera cima, la empresa de modelos pronto se convertiría en una empresa puramente de ingeniería; si no la tiene, la investigación, el producto y la organización seguirán enfrentando nuevos problemas. La montaña infinita significa que el espacio de problemas de largo plazo sigue existiendo.
\end{knowledgebox}

\subsection{Test-time scaling: el pensamiento largo y los Agent apuntan en la misma dirección}

La sección anterior trató de que «los problemas seguirán surgiendo»; esta entra en el primer juicio técnico. Yang Zhilin considera que tanto el aprendizaje por refuerzo basado en pensamiento largo como el aprendizaje por refuerzo de Agent apuntan hacia el test-time scaling. El test-time scaling consiste en invertir más cómputo en la etapa de inferencia/prueba: cadenas de razonamiento más largas, más llamadas a herramientas, más verificación y más iteraciones.

\begin{figure}[H]
\centering
\includegraphics[width=0.96\textwidth]{test-time-scaling-map.png}
\caption{Test-time Scaling: tanto el RL de pensamiento largo como el RL de Agent amplían el cómputo en tiempo de inferencia. Diagrama conceptual propio, elaborado a partir de la conversación de 00:03:38--00:24:58.}
\end{figure}

\begin{knowledgebox}{Lectura de la figura: el escalamiento en tiempo de prueba no consiste solo en generar más token}
El pensamiento largo hace que el modelo gaste más token en el razonamiento interno; los Agent permiten que el modelo obtenga retroalimentación externa a través de herramientas y del entorno. Ambos gastan más cómputo en la etapa de prueba; uno se inclina por el pensamiento interno y el otro por la acción externa.
\end{knowledgebox}

\begin{importantbox}{Fórmula intuitiva del test-time scaling}
\[
\text{Capability} \approx f(\text{model}, \text{test-time compute}, \text{feedback})
\]
Donde \(\text{model}\) es la capacidad del modelo base, \(\text{test-time compute}\) es la inversión adicional en la etapa de inferencia y \(\text{feedback}\) es la verificación, lo que devuelven las herramientas o la retroalimentación del entorno. Lo decisivo en un Agentic LLM es convertir estos dos últimos elementos en un sistema.
\end{importantbox}

\subsection{De L1 a L5 no necesariamente en serie}

La sección anterior descompuso el test-time scaling en pensamiento interno y acción externa; esta explica además que estas capacidades no necesariamente suben los peldaños de forma lineal. En la entrevista se menciona que de L1 a L5 no hay necesariamente una relación en serie: Claude logró muy buenos resultados en Agent aun cuando su reasoning no era necesariamente el más fuerte. Este juicio es importante: el desarrollo de las capacidades de los modelos no es una escalera fija; se puede apostar en paralelo por Reasoning, Agent, herramientas, producto y retroalimentación. Solo cuando el modelo participa en el proceso de desarrollo es posible desbloquear la verdadera etapa de Innovator.

\begin{figure}[H]
\centering
\includegraphics[width=0.96\textwidth]{l1-l5-agent-ladder.png}
\caption{De L1 a L5 no necesariamente en serie: se puede apostar en paralelo por Reasoning y Agent; Claude es un ejemplo típico. Diagrama conceptual propio, elaborado a partir de la conversación de 00:12:00--00:24:58.}
\end{figure}

\begin{warningbox}{No hay que tratar los niveles de capacidad como una escalera mecánica}
Si se entiende que de L1 a L5 hay que superar cada nivel en orden, se subestiman las rutas de producto y de sistema. La capacidad de Agent puede despuntar primero en ciertas tareas y, después, impulsar a su vez el entrenamiento del modelo y el ecosistema de herramientas.
\end{warningbox}

\begin{knowledgebox}{Explicación didáctica de las etapas L}
\begin{tabularx}{\textwidth}{p{0.16\textwidth}p{0.34\textwidth}X}
\toprule
Etapa & Intuición & Variables decisivas \\
\midrule
L1/L2 & Chat, preguntas y respuestas, reasoning básico & Base del modelo y alineación. \\
L3 & Puede usar herramientas para completar tareas bien definidas & Interfaces de herramientas, definición de tareas, retroalimentación. \\
L4 & Participa en el desarrollo, la innovación y la resolución de problemas complejos & Agent de largo alcance, verificación, capacidad de programación. \\
L5 & Sistemas con mayor autonomía & Seguridad, organización, permisos y retroalimentación continua. \\
\bottomrule
\end{tabularx}
\end{knowledgebox}

\subsection{Resumen de la sección}

«La montaña infinita» da a este episodio su tono de fondo: cada paso que avanza la capacidad del modelo abre nuevos problemas; el test-time scaling, por su parte, marca la dirección técnica: hacer que el modelo piense más, actúe más y obtenga más retroalimentación durante la inferencia.
\section{K2 es el Chogori: modelo base, eficiencia y capacidad Agentic}

Esta sección entra en K2. Yang Zhilin (杨植麟) divide el objetivo de K2 en varias partes: primero, que sea un modelo base muy bueno; segundo, aprovechar al máximo cada porción de datos, es decir, aumentar la token efficiency; tercero, que tenga buena capacidad Agentic y pueda pasar de ser un «cerebro en una cubeta» al uso de herramientas en múltiples turnos; cuarto, que el código abierto lleve esas capacidades al ecosistema.

\begin{figure}[H]
\centering
\includegraphics[width=0.96\textwidth]{k2-design-goals.png}
\caption{Objetivos de diseño de Kimi K2: modelo base, token efficiency, capacidad Agentic y código abierto deben cumplirse a la vez. Diagrama conceptual propio, elaborado a partir de la conversación entre 00:24:58 y 00:54:08.}
\end{figure}

\begin{knowledgebox}{Lectura de la figura: K2 no es una optimización aislada}
El modelo base determina los cimientos; la token efficiency determina cuánta capacidad puede surgir de los mismos datos; Muon y Rephrase son medios para mejorar la eficiencia del entrenamiento; la capacidad Agentic determina si el modelo puede usar herramientas, y el código abierto determina la difusión en el ecosistema.
\end{knowledgebox}

\subsection{Token efficiency: con los mismos datos, el cerebro crece más}

Después del muro de datos, la token efficiency se vuelve más importante. En palabras de Yang Zhilin, con la misma cantidad de datos se busca que el «cerebro» crezca más. K2 aplica Rephrase a los datos y también presta atención a líneas de trabajo como el optimizador Muon. Estas técnicas no buscan solo mejorar en los benchmarks, sino aumentar el rendimiento de entrenamiento de cada token a medida que los datos de alta calidad se vuelven más escasos.

\begin{figure}[H]
\centering
\includegraphics[width=0.96\textwidth]{token-efficiency-loop.png}
\caption{Ciclo de la Token Efficiency: la reescritura de datos, el optimizador y la estrategia de entrenamiento aumentan el rendimiento de cada token. Diagrama conceptual propio, elaborado a partir de la conversación entre 00:27:28 y 00:32:44.}
\end{figure}

\begin{knowledgebox}{Términos clave: conceptos relacionados con la eficiencia del entrenamiento}
\begin{tabularx}{\textwidth}{p{0.22\textwidth}p{0.32\textwidth}X}
\toprule
Concepto & Función & Riesgo \\
\midrule
Token efficiency & Obtener más capacidad con la misma cantidad de tokens & Es difícil de medir de forma aislada y suele confundirse con otras variables. \\
Rephrase & Reescribir los datos para aumentar la variedad de expresiones & Una reescritura de mala calidad introduce ruido. \\
Muon & Línea de trabajo en optimizadores que puede mejorar la eficiencia del entrenamiento & El entrenamiento puede volverse inestable; en la entrevista se menciona que «explota». \\
Data mixture & Proporción de los datos y currículo de entrenamiento & Una proporción equivocada puede dañar la capacidad o la generalización. \\
\bottomrule
\end{tabularx}
\end{knowledgebox}

\subsection{El reto de generalización de los modelos Agentic}

La sección anterior trató la eficiencia del entrenamiento de K2; esta pasa al problema central de la capacidad Agentic: la generalización. Para un modelo Agentic, el mayor reto es generalizar. Si las herramientas, las tareas y la retroalimentación del entorno de entrenamiento son demasiado limitadas, el modelo puede desempeñarse bien solo en escenarios conocidos; lo verdaderamente valioso es que, ante herramientas, tareas y entornos nuevos, siga siendo capaz de explorar en múltiples turnos, invocar herramientas y corregir su comportamiento.

\begin{figure}[H]
\centering
\includegraphics[width=0.94\textwidth]{agentic-generalization.png}
\caption{El reto de generalización de los modelos Agentic: saber usar herramientas no equivale a generalizar a tareas nuevas. Diagrama conceptual propio, elaborado a partir de la conversación entre 00:32:50 y 00:38:00.}
\end{figure}

\begin{importantbox}{Cómo juzgar la generalización Agentic}
Para saber si un modelo es verdaderamente Agentic no basta con ver si invoca una herramienta fija; hay que ver si, ante tareas desconocidas, puede entender la documentación de las herramientas, planificar acciones de varios pasos, corregirse a partir de la retroalimentación y transferir su capacidad a entornos nuevos.
\end{importantbox}

\begin{knowledgebox}{Tres niveles de prueba de la generalización Agentic}
\begin{tabularx}{\textwidth}{p{0.22\textwidth}p{0.34\textwidth}X}
\toprule
Nivel & Qué se evalúa & Modo de falla \\
\midrule
Generalización de herramientas & Si entiende API, comandos y documentación nuevos & Solo sabe usar las herramientas vistas en el entrenamiento. \\
Generalización de tareas & Si puede descomponer objetivos nuevos y dependencias de largo alcance & El plan se ve bien, pero no se puede ejecutar. \\
Generalización de la retroalimentación & Si puede corregirse después de recibir un error & Persiste en el error o repite la misma acción. \\
\bottomrule
\end{tabularx}
\end{knowledgebox}

\subsection{Resumen de la sección}

El núcleo de K2 es equilibrar varios objetivos: un modelo base sólido, eficiencia de entrenamiento, generalización Agentic y un ecosistema de código abierto deben cumplirse a la vez. Una debilidad en cualquiera de estos eslabones limita su paso de modelo a sistema.
\section{Plan de decisiones técnicas de K2: cómo se condicionan entre sí los cuatro objetivos}

La sección anterior explicó que K2 no es una optimización aislada; esta sección organiza esos objetivos en un plan de decisiones y responde a la pregunta «¿qué optimiza a la vez una empresa de modelos?». El modelo base, las capacidades agénticas, el ecosistema de código abierto y la comercialización no son consignas paralelas, sino cuatro objetivos que se condicionan entre sí. Cuanto más fuerte es el modelo base, más respaldo tienen los productos de Agent y la API; cuanto más completa es la apertura del código, más rápido se difunde el ecosistema, pero más difícil resulta la captura comercial; cuanto más fuertes son las capacidades agénticas, más importan las herramientas y la retroalimentación, y más complejos se vuelven el entrenamiento y la evaluación.

\begin{knowledgebox}{Matriz de los cuatro objetivos de K2}
\begin{tabularx}{\textwidth}{p{0.18\textwidth}p{0.30\textwidth}p{0.38\textwidth}}
\toprule
Objetivo & Qué se optimiza & Qué arrastra consigo \\
\midrule
Base Model & Capacidad general, capacidad de programación, capacidad de razonamiento & Calidad de los datos, token efficiency, optimizador y capacidad de cómputo. \\
Agentic & Uso de herramientas en varios turnos, retroalimentación del entorno, generalización entre tareas & RL, cómputo en tiempo de prueba, protocolos de herramientas y evaluación. \\
Open Source & Validación externa, ecosistema de desarrolladores, confianza & Captura comercial, costo de soporte y reutilización por parte de competidores. \\
Business & API, productos, pagos de empresas y desarrolladores & Límites del modelo, límites de las aplicaciones de Agent y estructura de costos. \\
\bottomrule
\end{tabularx}
\end{knowledgebox}

\begin{importantbox}{Por qué estos cuatro objetivos no pueden analizarse por separado}
Si solo hay un modelo base fuerte sin capacidades agénticas, el modelo sigue siendo como un «cerebro en una cubeta»; si solo hay productos de Agent con un modelo base débil, la generalización queda limitada; si hay código abierto pero no una vía comercial, el ecosistema difícilmente se sostiene; si la comercialización llega demasiado pronto, puede limitar la investigación y la difusión del código abierto.
\end{importantbox}

\subsection{El ciclo entre eficiencia de entrenamiento y capacidad del producto}

En la matriz de los cuatro objetivos, lo que más fácilmente se subestima es la eficiencia de entrenamiento; esta subsección explica por qué se transmite al producto. La token efficiency parece una métrica de entrenamiento, pero influye en la capacidad del producto. Una mayor token efficiency significa obtener un modelo base más fuerte con los mismos datos y la misma capacidad de cómputo; un modelo base más fuerte permite que el Agent resuelva mejor tareas complejas; las tareas complejas generan más retroalimentación y más datos de uso; y esa retroalimentación puede incorporarse a la siguiente ronda de entrenamiento y de iteración del producto. Si este ciclo se pone en marcha, la empresa de modelos ya no solo entrena modelos, sino que construye un volante de inercia de capacidades.

\begin{knowledgebox}{Volante de inercia de capacidades}
\begin{tabularx}{\textwidth}{p{0.20\textwidth}p{0.34\textwidth}X}
\toprule
Etapa & Función & Riesgo \\
\midrule
Eficiencia de entrenamiento & Reducir el costo de obtener capacidades & Inestabilidad del optimizador o de la estrategia de datos. \\
Producto de Agent & Llevar el modelo a tareas reales & Tasa de fallas, seguridad de las herramientas y confianza del usuario. \\
Datos de retroalimentación & Registrar éxitos y fallas en tareas reales & Mucho ruido, privacidad y sesgo de selección. \\
Siguiente ronda de entrenamiento & Corregir el modelo y la estrategia de herramientas con la retroalimentación & Una evaluación imprecisa amplifica los errores. \\
\bottomrule
\end{tabularx}
\end{knowledgebox}

\subsection{Resumen de la sección}

Las decisiones técnicas de K2 no buscan el óptimo de un solo punto, sino la optimización conjunta de cuatro objetivos. Solo al entender este plan se entiende por qué la entrevista va y viene entre la token efficiency, la generalización agéntica, el código abierto, la API y el modelo de negocio.
\section{Un sistema simple y complejo a la vez: código abierto, muro de datos y decisiones de arquitectura}

La sección anterior trató los objetivos de diseño de K2; esta examina las decisiones a nivel de sistema. ¿Por qué pasar de código cerrado a código abierto? ¿Por qué ya es un buen resultado que la multimodalidad no dañe el «cerebro»? ¿Por qué, cuando la Scaling Law se topa con el muro de datos, hay que reequilibrar los FLOPs, el RL, la retroalimentación y la eficiencia de la arquitectura? Estas preguntas muestran que una empresa de modelos no ajusta una sola perilla, sino que toma decisiones continuamente dentro de un sistema complejo.

\begin{figure}[H]
\centering
\includegraphics[width=0.94\textwidth]{open-source-strategy.png}
\caption{De código cerrado a código abierto: el grado de madurez del modelo y del producto, la difusión en el ecosistema y la vía comercial influyen en conjunto en la decisión de abrir el código. Diagrama conceptual propio, elaborado a partir de la conversación entre 00:54:08 y 01:05:00.}
\end{figure}

\begin{knowledgebox}{Lectura de la figura: abrir el código es una decisión técnica y también una decisión de ecosistema}
El código cerrado facilita controlar la experiencia del producto y la comercialización; el código abierto favorece la difusión en el ecosistema, la verificación por parte de los desarrolladores y la construcción de confianza. Que K2 sea de código abierto implica que debe demostrar la capacidad del modelo base y, al mismo tiempo, permitir que el ecosistema Agentic crezca a su alrededor.
\end{knowledgebox}

\subsection{Scaling Law, muro de datos y FLOPs}

La sección anterior trató el código abierto y el ecosistema; esta vuelve a las restricciones duras del entrenamiento de modelos. Yang Zhilin (杨植麟) considera que la Scaling Law se ha topado con el muro de datos, y que es un hecho objetivo. Cuando el crecimiento de los datos de alta calidad se desacelera, seguir haciendo scale de los datos ya no es tan sencillo como antes; por eso aumenta la importancia del FLOPs scaling, del RL, del feedback del entorno, del volante de datos y de la eficiencia de la arquitectura. Por ahora, el scaling basado en FLOPs sigue funcionando, pero el punto de equilibrio podría cambiar en el futuro.

\begin{figure}[H]
\centering
\includegraphics[width=0.96\textwidth]{scaling-data-wall.png}
\caption{Scaling y muro de datos: después del muro de datos, se reequilibran los FLOPs, el RL, la retroalimentación y la eficiencia de la arquitectura. Diagrama conceptual propio, elaborado a partir de la conversación entre 01:03:00 y 01:10:00.}
\end{figure}

\begin{warningbox}{El muro de datos no significa que el entrenamiento se detenga}
El muro de datos no quiere decir que ya no haya datos, sino que escasean los datos de alta calidad, con poco ruido y capaces de aportar una mejora marginal de capacidad. Las empresas de modelos todavía pueden avanzar mediante mejores proporciones de datos, datos sintéticos, RL, entornos de herramientas y test-time scaling.
\end{warningbox}

\begin{knowledgebox}{Cuatro caminos después del muro de datos}
\begin{tabularx}{\textwidth}{p{0.20\textwidth}p{0.34\textwidth}X}
\toprule
Ruta & Qué resuelve & Costo \\
\midrule
FLOPs scaling & Seguir bajando la loss con más cómputo & Costo alto; depende del suministro de hardware. \\
Data efficiency & Obtener más capacidad con los mismos datos & Requiere mejor filtrado, reescritura, proporciones y optimizadores. \\
RL/feedback & Crear nuevas señales de aprendizaje a partir de la retroalimentación del entorno & El ruido del reward y los problemas de seguridad son más difíciles. \\
Architecture & Reducir el costo de contexto largo e inferencia con nuevas estructuras & Puede introducir bias en la capacidad. \\
\bottomrule
\end{tabularx}
\end{knowledgebox}

\subsection{Linear Attention y el riesgo para el «coeficiente intelectual»}

Después del muro de datos, la eficiencia de la arquitectura se vuelve más atractiva; esta sección analiza el riesgo de ese atractivo. En la entrevista se menciona que muchas arquitecturas de long context afectan el «coeficiente intelectual»; una Linear Attention pura podría afectar la capacidad del modelo debido al bias de la arquitectura. Este juicio coincide con la revisión de arquitecturas de Attention del EP119: reducir el costo del contexto largo es importante, pero si el cambio de arquitectura sacrifica capacidad de razonamiento, de memoria o expresiva, puede afectar el techo del modelo.

\begin{figure}[H]
\centering
\includegraphics[width=0.94\textwidth]{linear-attention-iq-risk.png}
\caption{Arquitecturas de Long Context y riesgo para el coeficiente intelectual: la Linear Attention reduce costos, pero el bias puede afectar la capacidad del modelo. Diagrama conceptual propio, elaborado a partir de la conversación entre 01:17:17 y 01:20:00.}
\end{figure}

\begin{knowledgebox}{Términos clave: decisiones de arquitectura}
\begin{tabularx}{\textwidth}{p{0.22\textwidth}p{0.32\textwidth}X}
\toprule
Ruta & Beneficio & Riesgo \\
\midrule
Full Attention & Gran capacidad expresiva; experiencia consolidada & Costo alto en contexto largo; mucha presión sobre la KV cache. \\
Linear Attention & Costo bajo; adecuada para secuencias largas & El bias de la arquitectura puede afectar la capacidad. \\
Sparse Attention & Reduce el cómputo y conserva parte de la visión global & El mecanismo de selección y la eficiencia en hardware son complejos. \\
Hybrid & Punto intermedio entre eficiencia y capacidad & Proporciones, estabilidad del entrenamiento y evaluación más complejas. \\
\bottomrule
\end{tabularx}
\end{knowledgebox}

\subsection{La frontera entre las empresas de modelos base y las empresas de productos Agent}

Lo anterior trató las decisiones de arquitectura internas del modelo; esta sección lleva la pregunta a la frontera de la industria. A largo plazo, la frontera entre las empresas de modelos base y las empresas de aplicaciones Agent no es fija. Las empresas de modelos base controlan la capacidad del modelo, los datos de entrenamiento, la API y el ecosistema de plataforma; las empresas de aplicaciones Agent controlan los escenarios, los flujos de trabajo, el punto de acceso de los usuarios y los datos de retroalimentación. Quien logre convertir la capacidad del modelo en un ciclo cerrado de tareas reales ocupará la posición de mayor valor.

\begin{figure}[H]
\centering
\includegraphics[width=0.94\textwidth]{agent-product-boundary.png}
\caption{Empresas de modelos base vs. empresas de aplicaciones Agent: la frontera a largo plazo depende del modelo, los datos, las herramientas y el punto de acceso de los usuarios. Diagrama conceptual propio, elaborado a partir de la conversación entre 01:09:00 y 01:15:00.}
\end{figure}

\begin{importantbox}{Criterio sobre la frontera}
Si un producto Agent es solo una capa delgada de workflow, la empresa de modelos base absorberá esa frontera; si el producto Agent cuenta con escenarios sólidos, un fuerte retorno de datos y un punto de acceso de usuarios de uso frecuente, puede, por el contrario, definir las necesidades del modelo.
\end{importantbox}

\begin{knowledgebox}{Comparación de modelos de negocio: API, producto y ecosistema}
\begin{tabularx}{\textwidth}{p{0.20\textwidth}p{0.34\textwidth}X}
\toprule
Modelo & Ventaja & Riesgo \\
\midrule
API & Fácil de escalar; los desarrolladores pueden integrarse & Competencia en el precio por token; lejos del valor para el usuario. \\
Producto Agent & Más cerca de las tareas y del punto de acceso de los usuarios & Requiere producto, herramientas, retroalimentación y fiabilidad. \\
Ecosistema de código abierto & Difusión rápida; gran confianza y participación de desarrolladores & Más difícil capturar valor comercial; costo de soporte alto. \\
Producto de código cerrado & Experiencia controlable; ciclo comercial claro & Difusión en el ecosistema y verificación externa limitadas. \\
\bottomrule
\end{tabularx}
\end{knowledgebox}

El modelo de negocio tiene además un problema del lado de los costos: las tareas Agentic suelen consumir más cómputo en inferencia que una conversación común. Una pregunta y respuesta sencilla puede requerir una sola llamada al modelo; una tarea Agent compleja puede requerir planificación, búsqueda, llamadas a herramientas, verificación, reintentos y resumen. Esto amplía la brecha entre el «valor para el usuario» y el «costo de inferencia»: si la tarea es de alto valor, el costo del Agent puede absorberse; si la tarea es solo una conversación de bajo valor, el costo puede aplastar fácilmente el margen bruto.

\begin{knowledgebox}{Estructura de costos de un Agent}
\begin{tabularx}{\textwidth}{p{0.20\textwidth}p{0.34\textwidth}X}
\toprule
Rubro de costo & De dónde proviene & Cómo reducirlo \\
\midrule
Llamadas al modelo & Planificación, ejecución y reflexión en varias rondas & Enrutamiento a modelos pequeños, caché, división de tareas. \\
Llamadas a herramientas & Búsqueda, código, navegador, API & Mejores protocolos de herramientas y recuperación ante fallos. \\
Costo de verificación & Pruebas, revisiones, confirmación humana & Verificadores automáticos y clasificación por nivel de riesgo. \\
Costo de contexto & Historial largo, archivos, bases de código, páginas web & Compresión de memoria y estrategias de recuperación. \\
\bottomrule
\end{tabularx}
\end{knowledgebox}

\subsection{Resumen de la sección}

Las decisiones de sistema detrás de K2 abarcan la estrategia de código abierto, el muro de datos, el FLOPs scaling, la eficiencia de la arquitectura y la frontera del producto. No es una competencia técnica en un solo punto, sino una optimización con múltiples variables que la empresa de modelos realiza dentro de un sistema complejo.
\section{En su propia historia: gestión con RL, miedo y amplificador de la civilización}

Esta sección vuelve a la perspectiva del fundador. Tim dice que hay que gestionar al estilo de RL y no solo con SFT. La metáfora resulta interesante: SFT equivale a mostrarle al equipo las respuestas correctas, mientras que RL equivale a fijar objetivos y retroalimentación para que el equipo aprenda estrategias mediante la exploración. Sin embargo, el riesgo de gestionar con RL es que resulta fácil ser objeto de hack; es decir, el equipo puede optimizar los indicadores en lugar del objetivo real.

\begin{figure}[H]
\centering
\includegraphics[width=0.96\textwidth]{rl-management-sft-risk.png}
\caption{Gestionar con RL y no solo con SFT: los objetivos, la retroalimentación y la recompensa moldean al equipo, pero también son fáciles de someter a hack. Diagrama conceptual propio, elaborado a partir de la conversación de 01:25:05--01:34:00.}
\end{figure}

\begin{knowledgebox}{Lectura de la figura: por qué es útil la metáfora de la gestión}
La gestión tipo SFT pone el acento en la demostración y en las respuestas correctas; la gestión tipo RL lo pone en los objetivos, la retroalimentación y la exploración. Tanto el entrenamiento de modelos como la gestión de organizaciones necesitan feedback, pero si el feedback está mal diseñado, termina siendo objeto de hack.
\end{knowledgebox}

\subsection{Complejidad e historia}

La subsección anterior explicó que la gestión con RL moldea al equipo; esta vuelve a cómo entiende el fundador esa complejidad. Yang Zhilin (杨植麟) dice que gran parte de la complejidad se añade de forma artificial y que, en realidad, las cosas no son tan complejas; aun así, el fundador necesita sentir, dentro de su propia historia, qué clase de persona es y por qué hace lo que hace. Esto no es huir de la técnica, sino reconocer que una empresa de modelos es la superposición de un sistema técnico y un sistema narrativo: la ruta técnica exige criterio, la organización necesita sentido y el fundador necesita seguir actuando en medio de la incertidumbre.

\begin{importantbox}{La IA es un amplificador de la civilización humana}
En la entrevista se menciona que la respuesta de Kimi a «¿qué es la IA?» es: la IA es un amplificador de la civilización humana. Esta definición eleva el modelo de la escala de herramienta a la escala de civilización: amplifica el conocimiento, la productividad y la capacidad organizativa, pero también amplifica los miedos, los malentendidos y las controversias propias de los estados intermedios.
\end{importantbox}

\begin{knowledgebox}{Tres tipos de complejidad en la historia del fundador}
\begin{tabularx}{\textwidth}{p{0.20\textwidth}p{0.34\textwidth}X}
\toprule
Complejidad & Origen & Forma de afrontarla \\
\midrule
Complejidad técnica & El modelo, los datos, la arquitectura y el producto cambian al mismo tiempo & Volver a problemas verificables y al siguiente experimento. \\
Complejidad organizativa & El equipo, la retroalimentación, la opinión pública y los indicadores de gestión se influyen entre sí & Corregir con retroalimentación, pero evitar que el reward sea objeto de hack. \\
Complejidad narrativa & Los estados intermedios pueden ser malinterpretados o criticados desde fuera & Mantener la dirección de la acción dentro de la propia historia. \\
\bottomrule
\end{tabularx}
\end{knowledgebox}

\subsection{El miedo y el paso actual}

Ante las tormentas de opinión pública y los altibajos del emprendimiento, Yang Zhilin admite que sin duda hay miedo, pero que lo más importante es concentrarse en lo que se puede hacer en el paso actual. Esta actitud se corresponde con la idea de «la montaña infinita»: si los problemas son inevitables, el miedo también lo es; pero los problemas pueden resolverse, de modo que la acción debe volver al siguiente paso.

\begin{warningbox}{Los estados intermedios siempre serán criticados}
Todo sistema complejo es imperfecto en sus estados intermedios y puede convertirse en blanco de críticas. Esto vale en especial para las empresas de modelos: las capacidades, el producto, el código abierto, la comercialización, la organización y la opinión pública se encuentran en un equilibrio dinámico. No se puede dejar de resolver problemas solo porque un estado intermedio reciba críticas.
\end{warningbox}

\subsection{Resumen de la sección}

La última sección devuelve el juicio técnico a la historia del fundador. K2 no es solo el lanzamiento de un modelo, sino también una narrativa sobre sí mismo que trata de problemas infinitos, retroalimentación organizativa, gestión del miedo y la IA como amplificador de la civilización.
\section{Síntesis y ampliación}

Esta sección condensa el episodio completo en cinco conclusiones. Primera: lo esencial de un Agentic LLM es pasar de ser un «cerebro en una cubeta» a un actor digital capaz de usar herramientas. Segunda: el test-time scaling abarca tanto el razonamiento prolongado como la retroalimentación de las herramientas que usa el Agent. Tercera: K2 reúne en un mismo diseño el modelo base, la token efficiency, Muon/Rephrase, la generalización Agentic y el ecosistema de código abierto. Cuarta: después del muro de datos, los FLOPs, el RL, la retroalimentación y la eficiencia de la arquitectura volverán a equilibrarse, y las vías de bajo costo, como Linear Attention, deben vigilar el bias en las capacidades. Quinta: una empresa de modelos no solo debe resolver problemas técnicos, sino también los del código abierto, el modelo de negocio, la retroalimentación organizacional y la mentalidad del fundador.

\begin{knowledgebox}{Lista de juicios principales}
\begin{tabularx}{\textwidth}{p{0.24\textwidth}p{0.34\textwidth}X}
\toprule
Juicio & Por qué se sostiene & Qué falta observar \\
\midrule
Lo Agentic es la siguiente línea principal & Los modelos ya pueden escribir código, usar herramientas y ejecutar tareas de largo alcance & Si la generalización y la confiabilidad pueden seguir mejorando. \\
La token efficiency cobra importancia & Los datos de alta calidad se vuelven escasos en el margen & Si la reescritura de datos y el optimizador son estables. \\
El código abierto tiene valor estratégico & La validación del ecosistema y la difusión entre desarrolladores son muy fuertes & La captura de valor comercial y el costo de mantenimiento. \\
La eficiencia de la arquitectura no puede sacrificar la inteligencia & El costo del contexto largo tiene que bajar & El techo de capacidades de Linear/Sparse/Hybrid. \\
Las fronteras comerciales se redibujarán & La API, los productos de Agent y el ecosistema se condicionan entre sí & Quién controla el punto de entrada de los usuarios y los datos de retroalimentación. \\
\bottomrule
\end{tabularx}
\end{knowledgebox}

\begin{importantbox}{Ubicar el EP113 en la serie de IA de Zhang Xiaojun (张小珺)}
El EP119 trata la arqueología de las arquitecturas de Attention, el EP115 la teoría de la segunda mitad de los Agent y el EP116 el Agentic Model para empresas; el EP113, en cambio, ofrece la perspectiva de una empresa de modelos: cómo K2 combina el modelo base, las capacidades Agentic y el ecosistema de código abierto, y cómo busca nuevas vías de scaling después del muro de datos.
\end{importantbox}

\subsection{Puntos principales}

\begin{enumerate}
\item El valor de K2 está en la combinación de las capacidades del modelo base, la programación y el uso de herramientas, la generalización Agentic y el ecosistema de código abierto.
\item El test-time scaling es una abstracción común al razonamiento prolongado y al uso de herramientas por parte del Agent.
\item La token efficiency es uno de los indicadores centrales del entrenamiento en la era del muro de datos.
\item En las arquitecturas de contexto largo, como Linear Attention, hay que considerar a la vez la eficiencia y el riesgo para la «inteligencia».
\item Las fronteras entre el código abierto, la API, los productos de Agent y el modelo base son una cuestión central para la comercialización de una empresa de modelos.
\end{enumerate}

\subsection{Relación con los episodios anteriores y posteriores}

\begin{tabularx}{\textwidth}{p{0.18\textwidth}p{0.34\textwidth}X}
\toprule
Episodio & Tema & Conexión con el EP113 \\
\midrule
EP119 & Arquitecturas de atención de Kimi Linear, DeepSeek y MiniMax & Explica las decisiones de arquitectura de Attention y de contexto largo detrás de K2. \\
EP115 & La segunda mitad de los Agent, reward, interface & Ofrece el marco teórico del Agentic LLM. \\
EP116 & Agentic Model para empresas & Muestra la versión aplicada de las capacidades Agentic en escenarios ToB con datos privados. \\
EP118 & VLA, Agent OS, puntos de entrada al mundo físico & Lleva la idea de que «al modelo le crecen manos» del mundo digital al mundo físico. \\
\bottomrule
\end{tabularx}

\subsection{Preguntas abiertas}

Las secciones anteriores presentaron conclusiones; esta reúne lo que de verdad sigue sin resolverse. No son un adorno de cierre, sino preguntas que, después de K2, tanto las empresas de modelos como las de productos de Agent deberán seguir respondiendo.

\begin{enumerate}
\item Después del muro de datos, ¿cómo se redistribuirá el presupuesto entre el FLOPs scaling, el RL feedback y la eficiencia de la arquitectura?
\item Entre Linear Attention, Sparse Attention y Hybrid Attention, ¿qué vía logrará imponerse en el equilibrio entre capacidad y costo?
\item ¿Cómo puede un modelo de código abierto encontrar un equilibrio entre la difusión en el ecosistema y la captura de valor comercial?
\item ¿Cómo debe evaluarse la generalización de un Agentic LLM, en particular ante herramientas y tareas desconocidas?
\item ¿Absorberán las empresas de modelos base a las aplicaciones de Agent, o serán las aplicaciones de Agent las que definan lo que se exige de los modelos base?
\end{enumerate}

\begin{warningbox}{El valor de las preguntas abiertas}
Estas preguntas no son un apéndice, sino el núcleo de las decisiones reales de una empresa de modelos. La importancia de K2 radica precisamente en que expone todas estas preguntas a la vez: ninguno de los aspectos (modelo, datos, herramientas, ecosistema, negocio y organización) puede resolverse por separado.
\end{warningbox}

\begin{knowledgebox}{Tabla anexa de correspondencia de términos}
\begin{tabularx}{\textwidth}{p{0.22\textwidth}p{0.34\textwidth}X}
\toprule
Término & Explicación para recordar & Dónde aparece en este episodio \\
\midrule
Cerebro en una cubeta & Modelo que solo tiene cognición, sin interfaz para actuar & Introducción y línea principal del Agentic LLM. \\
Constructor universal & Actor capaz de cambiar el mundo mediante herramientas & Capacidades de Agent y analogía con los seres humanos. \\
Muro de datos & Los datos de entrenamiento de alta calidad se vuelven escasos en el margen & Discusión sobre el scaling y el volante de inercia de datos. \\
Riesgo para la inteligencia & Una arquitectura que reduce costos puede dañar las capacidades del modelo & Decisiones sobre Linear Attention. \\
Gestión tipo RL & Moldear la exploración del equipo con objetivos y retroalimentación & Metáfora organizacional del fundador. \\
\bottomrule
\end{tabularx}
\end{knowledgebox}

\subsection{Lecturas adicionales}

\begin{itemize}
\item Si interesan las decisiones de arquitectura de Attention, conviene contrastarlas con la revisión de las arquitecturas de atención de Kimi Linear / MiniMax M2 / DeepSeek del EP119.
\item Si interesa la teoría de los Agent, conviene contrastarla con la entrevista a Yao Shunyu (姚顺雨) del EP115, para entender el reward, el multi-agent y la interface.
\item Si interesa la implementación en empresas, conviene contrastarla con la entrevista a Wu Minghui (吴明辉) del EP116, para observar la forma que adopta el Agentic Model en escenarios ToB con datos privados.
\end{itemize}

\end{document}
